\chapter{Заключение}
\label{sec:Chapter6} \index{Chapter6}

Целью работы являлось повышение качества и интерпретируемости данных
дистанционного зондирования путём трансфера радиолокационных изображений в
оптический диапазон средствами генеративных нейронных сетей. Цель достигнута,
задачи, поставленные во введении, решены. Ниже кратко изложены полученные
результаты.

\begin{enumerate}[label*=\arabic*.,leftmargin=2.0em,itemsep=0.3em,topsep=0.3em]
  \item Систематизированы методы трансфера SAR\,$\rightarrow$\,optical~--- от парных
        и непарных GAN (Pix2Pix, CycleGAN) и вейвлет-ориентированных архитектур
        (CFRWD-GAN) до гибридных CNN--Transformer и диффузионных моделей; для каждого
        семейства выделены сильные стороны и ограничения, существенные именно для
        радиолокационных данных.
  \item Подготовлены два набора парных данных. Основной~--- подмножество SEN1-2 из
        пяти репрезентативных сцен (11\,315 пар); вспомогательный QXS-SAROPT
        использован для проверки переносимости. Реализованный конвейер предобработки
        включает предварительное выравнивание пар в градиентной области (каскад
        phaseCorrelate\,$\rightarrow$\,ECC), снабжённое предохранителем «без регресса»;
        это смягчает остаточную геометрическую разрегистрацию, не ухудшая исходно
        согласованные пары.
  \item Адаптированы под специфику SAR архитектуры Pix2Pix и CFRWD-GAN, а также
        разработан собственный генератор \textbf{LLW-Former}~--- одношаговый и
        вейвлет-обусловленный. Стандартный stride-4 patch-embed в нём заменён
        двухуровневым Хаар-стемом, а выходная голова восстанавливает изображение
        точным обратным Хаар-преобразованием. Такой частотно-симметричный контур
        согласует базис генератора с функцией потерь и допускает их вычисление прямо
        в вейвлет-области. Функция потерь скомпонована из пяти слагаемых~--- пиксельного
        в Хаар-базисе, MS-SSIM, LPIPS, Focal Frequency и PatchNCE; ядром предложенного
        решения служит высокочастотный дискриминатор HF-D, навязывающий когерентную
        высокочастотную детализацию.
  \item Экспериментальная часть охватила абляцию всего семейства, серию
        контролируемых сравнений и сопоставление с опубликованными методами. Лучшая
        конфигурация, LLW-Former Base, даёт PSNR 18,54~дБ, SSIM 0,432, LPIPS 0,241 и
        FID 58,5; облегчённая Tiny ($\sim$29~млн параметров)~--- PSNR 17,28~дБ,
        SSIM 0,369, FID 73,0. По перцептивным метрикам отрыв от опубликованных GAN- и
        диффузионных подходов на SEN1-2 заметен (FID 58--73 против 116--265), тогда как
        по пиксельным модель держится вровень с современными диффузионными методами,
        оставаясь одношаговой и кратно более быстрой на инференсе.
  \item Показана практическая применимость метода: лёгкая (29--88~млн параметров) и
        одношаговая модель с малой стоимостью вывода годится для оперативного
        мониторинга, аугментации обучающих выборок и приведения радиолокационных
        снимков к привычному для интерпретатора виду.
\end{enumerate}

\medskip
\noindent\textbf{Научная новизна} работы складывается из трёх результатов.
Первый~--- высокочастотный дискриминатор HF-D, действующий на высокочастотном остатке
оптического изображения. Он повышает резкость и перцептивное качество (в
контролируемом сравнении LPIPS снижается на 8,8\,\%, FID~--- примерно на 11\,\%) и при
этом не удорожает инференс, поскольку задействован только на обучении. Второй~---
частотно-симметричная архитектура с замкнутым Хаар-контуром «вход\,--\,выход»,
приводящая пространственно-частотный базис генератора в соответствие с функцией
потерь. Третий результат носит отрицательный характер, но самостоятельно значим:
установлено, что плотная обучаемая регистрация пар SAR\,/\,оптика неосуществима в
принципе~--- модальный разрыв таков, что локальная кросс-модальная корреляция близка к
нулю; вдобавок показано, что инвариантные к сдвигу статистические потери лишь ранжируют
резкость, но не порождают её. Именно эти наблюдения и оправдывают опору на
адверсариальный механизм как на единственный работающий.

\medskip
\noindent\textbf{Направления дальнейшей работы.} Естественное продолжение~---
наращивание масштаба обучающих данных: переход от пяти сцен к девяти, а затем и к
полному датасету. Отдельно стоит проверить переносимость архитектуры на QXS-SAROPT и
распространить HF-D на конфигурацию Base, в том числе в связке с генератором повышенной
ёмкости. Представляет интерес и чувствительность результата к выбору вейвлет-базиса
(Хаар\,$\rightarrow$\,Daubechies). Наконец, для слаботекстурных областей, прежде всего
водных, где отображение SAR\,$\rightarrow$\,optical по своей природе многозначно,
оправдано введение явной модели алеаторической неопределённости.

\newpage